\begin{proof}

\begin{enumerate}[label={$\langle1\rangle$\arabic*.},leftmargin=*]
\item A transverse cusp has exactly one normalization
branch. The map from the horizontal section to the
normalization of its image consequently has degree
one. Both curves are smooth, so this finite map
identifies $C\simeq B$ with $D^\nu$. Its genus is
positive. Statement \ref{cdp:st:A.5.14}, including its actual
delta-one lci comparison, and \ref{cdp:st:A.5.5} therefore make
$D$ itself smooth elliptic and its cusp transversal
constant at every point. Thus $C\to D$ is an
isomorphism, without any assumption about the
ambient Picard pullback.

\item For the exact local normalization $x=u^3,y=u^2$,
let $C_T\subset T$ be the image of $C$. The
normalization is smooth near this curve. Its normal
line $P=N_{C_T/T}$ has $P^2=R$ and $P^3=Q$
in the degree-zero radical and Hessian quotient
lines of \ref{cdp:st:C2d}. Hence $\deg P=0$. The normalization
is smooth near $C$, and no other conductor curve
meets it. Factor $\pi$ through its minimal
resolution, descending the ruling across its
vertical exceptional blowdowns as in \ref{cdp:st:A.5.18}.
On that minimal resolution the section is isomorphic
to $C_T$ in its unchanged smooth neighbourhood. It
therefore has self-intersection zero and is disjoint from
the effective vertical conductor. The canonical
comparison is the actual line-bundle identity
$K+2C+V=(\eta\pi)^*\omega_S$, of zero rational
Chern class.

\item Apply (A.T.22) to each ruling blowdown. Its
exceptional coefficient is exactly the multiplicity
at its centre minus one. No ruling $(-1)$-curve
can meet $C$, since the intersection with $K+2C+V$
would then be positive. On the minimal ruled
bundle the vertical conductor is zero because it
is disjoint from a section. A first reverse blowup
away from that section would have coefficient
$-1$, so there were no ruling blowups. All
resolution-exceptional curves are vertical, and
the square-zero fibres of the resulting bundle
cannot be exceptional. This identifies the minimal
resolution with $T$ itself and proves the asserted
smoothness and minimality.

\item Finally the normalization is an isomorphism off
$C$ and $D$. The exact Euler comparison is
$e(S)=e(T)-e(C)+e(D)$; all three terms vanish,
since $T$ is ruled over an elliptic curve and
$C,D$ are elliptic. Thus $e(S)=0$, also for the
underlying topological space of the full fibre $mS$.

\end{enumerate}
\end{proof}
